
\def \Subject {آشنایی با محاسبات انجام شده در فرآیند آموزش }
\def \Session { یک}
% \setcounter{chapter}{\Session-1}
\renewcommand{\chaptername}{سوال}
\chapter{\Subject}
این سوال به روابط ریاضی مورد استفاده در یک شبکه عصبی پیچشی می‌پردازد.
\section{انجام یک مرحله آموزش کامل}
\subsection{مقادیر اولیه}
با استفاده از کد داده شده، اطلاعات شبکه مطابق زیر است:
\begin{latin}
	~\\
	$
	X_1 = \begin{bmatrix}
		-2 & 0 & -1 \\
		-2 & -2 & 0 \\
		0 & -2 & -1
	\end{bmatrix}
	\quad 
	X_2 = \begin{bmatrix}
		0 & 0 & 0 \\
		0 & 0 & 2 \\
		0 & 0 & 0
	\end{bmatrix}
	\\\\\\
	W_{conv1} = \begin{bmatrix}
		-1 & 0 \\
		1 & 2
	\end{bmatrix}
	\quad 
	W_{conv2} = \begin{bmatrix}
		1 & 1 \\
		0 & 2
	\end{bmatrix}
	\quad
	W_{linear} =
	\begin{bmatrix}
		2 &
		-2 &
		-4 &
		3
	\end{bmatrix}
	\\
	bias_{conv} = -1 
	\quad 
	bias_{linear} = 1
	$
\end{latin}

\begin{figure}[H]
	\centering
	\includegraphics[width=\textwidth]{images/q1_arch.jpg}
	\caption{
		معماری مدل
	}
	\label{fig:q2_two_layer}
\end{figure}
\subsection{انتشار رو به جلو}
برای بخش 
\lr{padding}
داریم:
\begin{latin}
	~\\
	$
		X_{{padded}_1} = \begin{bmatrix}
-2 & -2 & -2 & 0 & -2 \\
0 & -2 & 0 & -1 & 0 \\
-2 & -2 & -2 & 0 & -2 \\
-2 & 0 & -2 & -1 & -2 \\
-2 & -2 & -2 & 0 & -2 
	\end{bmatrix}
	\quad 
	X_{{padded}_2} = \begin{bmatrix}
		0 & 0 & 0 & 2& 0 \\
		0 & 0 & 0 & 0& 0 \\
		0 & 0 & 0 & 2& 0 \\
		0 & 0 & 0 & 0& 0 \\
		0 & 0 & 0 & 2& 0 
	\end{bmatrix}
	$
\end{latin}
سپس برای 
\lr{convolution }
داریم:
\begin{latin}
	~\\
	$
	z_1 = X_{{padded}_1} * W_1 = \begin{bmatrix}
		-2 & 0 & 0 & -1 \\
		-6 & -4 & -2 & -3 \\
		0 & -2 & -2  & -5 \\
		-4 & -6 & 0 & -3
	\end{bmatrix}
	\quad 
	z_2 = \begin{bmatrix}
		0 & 0 & 2 & 2 \\
		0 & 0 & 4 & 0 \\
		0 & 0 &  2 & 2 \\
		0 & 0 & 4 & 0
	\end{bmatrix}
	\\
	\Rightarrow
	z = z_1 + z_2 + bias = \begin{bmatrix}
		-3 & -1 & 1 & 0  \\
		-7 & -5 & 1 & -4 \\
		-1 & -3 & -1  & -4 \\
		-5 & -7 & 3 & -4
	\end{bmatrix}
	$
\end{latin}
بعد از این مرحله باید 
\lr{max pooling}
انجام شود:
\begin{latin}
	$
	z_{pooled} = \begin{bmatrix}
		-1 & 1 \\
		-1 & 3
	\end{bmatrix}
$
\end{latin}
سپس اعمال تابع 
\lr{leaky ReLU}
با شیب 
$0.1$

و سپس 
\lr{flat}
کردن را داریم:
\begin{latin}
	~\\
$
	a_1 = \begin{bmatrix}
		-0.1 & 1 \\
		-0.1 & 3
	\end{bmatrix}
	\Rightarrow
	a_{1_{flat}} = \begin{bmatrix}
		-0.1 \\
		1 \\
		-0.1 \\
		3
	\end{bmatrix}
	$
\end{latin}
حال باید لایه خطی نهایی را اعمال کرد و مقدار بایاس را به آن اضافه نمود:
\begin{latin}
	~\\
	$
z_2 = 		\begin{bmatrix}
	2 &
	-2 &
	-4 &
	3
\end{bmatrix} * \begin{bmatrix}
-0.1 \\
1 \\
-0.1 \\
3
\end{bmatrix}
= 7.2
\qquad 
a_2 = z_2 + bias = 8.2
	$
\end{latin}
پس یعنی خروجی شبکه برای این ورودی مقدار $8.2$ است.
\subsection{انتشار رو به عقب}

اول باید مقدار خطا حساب شود. با توجه به اینکه در سوال آمده است برای خطا از میانگین مجذور استفاده شود، مقدار برچسب برابر 
$-1$
 و ضریب یادگیری برابر 
 $0.1$
 است، داریم:
 
 \begin{latin}
 	~\\
 	$
 	error_{mse} = \cfrac{1}{n} \sum\limits_{i = 1}^n (y_i - \hat{y}_i)^2 
 	= 
 	9.2^2 = 84.64
 	$
 \end{latin}
 حال باید با استفاده از قاعده زنجیری، مقدار گرادیان خطا نسبت به هر پارامتر را یافت. (درموقع محاسبه مشتق و گذر از لایه 
 \lr{pooling}
 یک وضعیت تساوی پیش آمد. با مشورت با هوش مصنوعی، مقدار گرادیان فقط به اولین بیشینه اختصاص داده شد.
 )
 (برای کانولوشن از نماد $*$ استفاده شده است.)
 \begin{latin}
 	~\\
 	$
 	\cfrac{\partial E }{\partial a_2} = 2(8.2 - (-1)) = 18.4
 	\cfrac{\partial E}{\partial bias_{linear}} = \cfrac{\partial E}{\partial a_2} * \cfrac{\partial a_2}{\partial bias_{linear}}  = 18.4
 	\\
 	\cfrac{\partial E}{\partial W_{linear}} =
 	\cfrac{\partial E}{\partial a_2} * \cfrac{\partial a_2}{\partial W_{linear}}
 	= 
 	18.4 * a_{1_{flat}}^T = \begin{bmatrix}
 		-1.84 &
 		18.4 &
 		-1.84 &
 		55.2
 	\end{bmatrix}
 	\\
 	\cfrac{\partial E}{\partial a_{1_{flat}}} = 
 	\cfrac{\partial E}{\partial a_2} *
 	\cfrac{\partial a_2}{\partial z_2} *
 	\cfrac{\partial z_2}{\partial a_{1_{flat}}} = 
 	18.4 * 1 * 	
 	W_{linear}^T = 	\begin{bmatrix}
 		36.8 \\
 		-36.8 \\
 		-73.6 \\
 		55.2
 	\end{bmatrix}
 	\Rightarrow 
 	\cfrac{\partial E}{\partial a} = \begin{bmatrix}
 		36.8 &
-36.8 \\
-73.6 &
55.2
 	\end{bmatrix}
 	\\
 	\Rightarrow
 	\cfrac{\partial E}{\partial z_{pooled}} = \begin{bmatrix}
 		3.68 & -36.8 \\
 		-7.36 & 55.2
 	\end{bmatrix}
 	\\
 	\cfrac{\partial E}{\partial z} = 
 	\cfrac{\partial E}{\partial z_{pooled}} *
 	\cfrac{\partial z_{pooled}}{\partial z} = \begin{bmatrix}
 		0 & 3.68 & -36.8 & 0 \\
 		0 & 0 & 0 & 0 \\
 		-7.36 & 0 & 0 & 0 \\
 		0 & 0 & 55.2 & 0 
 	\end{bmatrix} 
 	\\
 	\cfrac{\partial E}{\partial bias_{conv}} = 
 	\cfrac{\partial E}{\partial z} * 
 	\cfrac{\partial z}{\partial bias_{conv}} = 
 	sum
 	 \begin{bmatrix}
 		0 & 3.68 & -36.8 & 0 \\
 		0 & 0 & 0 & 0 \\
 		-7.36 & 0 & 0 & 0 \\
 		0 & 0 & 55.2 & 0 
 	\end{bmatrix} 
 	= 14.72
 	\\
 	 	\cfrac{\partial E}{\partial z_1} = 
 	\cfrac{\partial E}{\partial z} * 
 	\cfrac{\partial z}{\partial z_1} 
 	= 
 	\cfrac{\partial z}{\partial z} 
 	\Rightarrow 
 	\cfrac{\partial E}{\partial z_1}  = 
 	\cfrac{\partial E}{\partial z_2} = 
 	\cfrac{\partial E}{\partial z} = 
 	 \begin{bmatrix}
 		0 & 3.68 & -36.8 & 0 \\
 		0 & 0 & 0 & 0 \\
 		-7.36 & 0 & 0 & 0 \\
 		0 & 0 & 55.2 & 0 
 	\end{bmatrix}   
 	\\
 	\cfrac{\partial E}{\partial W_{conv1}} = 
 	\cfrac{\partial E}{\partial z_1} * 
 	\cfrac{\partial z_1}{\partial W_{conv1}} = \cfrac{\partial E}{\partial z_1} *  X_{{padded}_1} 
 	 =  \begin{bmatrix}
 	 	-29.44 & -47.84 \\
 	 	-103.04 & 36.8
 	 \end{bmatrix}
 	 \\
 	\cfrac{\partial E}{\partial W_{conv2}} = 
 	\cfrac{\partial E}{\partial z_2} * 
 	\cfrac{\partial z_2}{\partial W_{conv2}} = \cfrac{\partial E}{\partial z_2} *  X_{{padded}_2} 
 	 =  \begin{bmatrix}
 	 	0 & -73.6 \\
 	 	0 & 110.4
 	 \end{bmatrix}
 	$
 \end{latin}
 حال با توجه به اینکه ضریب یادگیری برابر 
 $0.1$
 است، باید مقدار گرادیان‌ها را در 
 $0.1$
 ضرب کرد و با مقدار قبلی پارامتر جمع کرد تا به مقدار جدید رسید. پس داریم:
 
 \begin{latin}
 	~\\
 	$
 	bias_{linear} = 1 - 0.1 * 18.4 = 1 - 1.84 = -0.84 \\
 	W_{linear} = \begin{bmatrix}
 		2 & -2  & -4 & 3
 	\end{bmatrix} 
 	- 
 	0.1 \begin{bmatrix}
 		 		-1.84 &
 		18.4 &
 		-1.84 &
 		55.2
 	\end{bmatrix}
 	= 
 	 \begin{bmatrix}
 		2.184 & 
 		-3.84 &
 		-3.816 &
 		-2.52
 	\end{bmatrix}
 	\\
 	bias_{conv} = -1 - 0.1 * 14.72 = -2.472
 	\\
 	W_{conv1} = \begin{bmatrix}
 		-1 & 0 \\
 		1 & 2
 	\end{bmatrix}
 	- 
 	0.1 
 	 \begin{bmatrix}
 		-29.44 & -47.84 \\
 		-103.04 & 36.8
 	\end{bmatrix}
 	= 
 	 \begin{bmatrix}
 		1.944 & 4.784 \\
 		11.304 & -1.68
 	\end{bmatrix}
 	\\
 	W_{conv2} = 
 	 \begin{bmatrix}
 		1 & 1 \\
 		0 & 2
 	\end{bmatrix}
 	- 
 	0.1
 	  \begin{bmatrix}
 		0 & -73.6 \\
 		0 & 110.4
 	\end{bmatrix}
 	=
 	  \begin{bmatrix}
 		1 & 8.36 \\
 		0 & -9.04
 	\end{bmatrix}
 	$
 \end{latin}
 \section{مقایسه انواع \lr{pooling}}
 اگر از 
 \lr{average pooling}
 استفاده می‌شد، 
 در انتشار رو به جلو، هنگام بدست آوردن 
 $z_{pooled}$
 از روی 
 $z$
 به جای اینکه بیشینه درنظر گرفته شود، میانگین حساب می‌شود و در انتشار رو به عقب، به جای این‌که یک ماسک بولین
 \LTRfootnote{boolean}
 درنظر گرفته شود که مقادیر بیشینه در یک مربع دو در دو در آن یک هستند،‌مقدار گرادیان بر تعداد درایه‌های پنجره (در اینجا ۴) تقسیم 
 شده و مشتق 
 $z$
 بدست می‌آید.
 \section{تغییر تعداد کانال}
 در اینجا در لایه کانولوشن عملا فقط یک ویژگی حساب می‌شود و لایه فقط یک هسته 
 \LTRfootnote{kernel}
  دارد. برای تعداد کانال خروجی بیشتر، یعنی لایه کانولوشن باید ویژگی‌های بیشتری را بیابد، پس باید چندین هسته داشته باشد. هر هسته به اندازه کانال‌های ورودی * ابعاد فیلتر باید پارامتر داشته باشد.
 
 برای مثال، اگر بخواهیم خروجی لایه کانولوشن دو کاناله باشد، باید ۲ * ۲ * ۲ پارامتر جدید برای پارامترهای کانولوشن و ۱ پارامتر جدید برای بایاس هسته جدید معرفی کرد. هر هسته بین پارامترهای خود و کانال‌های ورودی عمل کانولوشن را انجام می‌دهد و سپس خروجی کانال‌های مختلف ورودی  و مقدار  \lr{bias}  را باهم جمع می‌کند. چون دو هسته موجود است، خروجی لایه یک خروجی دو کاناله با ابعاد ۴ در ۴ است. 
 